\section{The finite coupling renormalization}
\label{sec:renormalization}

The comparison in Section~\ref{sec:comparison} supplies parameters
$b=\beta+O(1)$ and $r=1+O(\beta^{-1})$. An additive constant in
$\beta-b$ changes the mass prefactor, so the boundedness of this difference
is not enough. We determine its limit by measuring the response to a small
quaternion twist. The calculation has two parts: the finite Taylor maps
show that the responses of the matched models agree through their constant
terms, and a direct Gaussian calculation evaluates their difference in the
microscopic variables.

\subsection{A helicity probe of the matched trajectories}

Let $e$ be a unit imaginary quaternion. On a square with $n$ initial cells
in each direction, impose the boundary condition
\[
 q_{x+ne_i}=q_x\exp(\zeta e),\qquad i\in\{x,t\},
\]
with periodic boundary conditions in the other direction. The twist acts
only on the quaternion field. In the Brownian model the physical horizontal
circumference is $nr/b$; the square terminology refers to the initial cell
coordinates. Write $Z_i^M(n;\zeta)$ for the resulting partition function in
model $M\in\{F,B\}$, and define its helicity by
\[
 h_i^M(n)=-\left.\frac{\partial^2}{\partial\zeta^2}
                       \log Z_i^M(n;\zeta)\right|_{\zeta=0}.
\]
For $F$ we use the flat-lift model without its phase. Left and right
quaternion twists have the same helicity: inversion of the quaternion
field interchanges the two conventions and preserves the microscopic
weights.

\begin{proposition}[Helicity matching]
\label{prop:matching-helicity}
Fix the blocking factor $L$ and the admission constants as in
Sections~\ref{sec:rg}--\ref{sec:comparison}. For every sufficiently large
$\beta$, choose any matched trajectory satisfying \eqref{eq:R15}, with its
parameters $b,r,c_x,c_t$. For $i=x,t$,
\begin{equation}
 \lim_{m\to\infty}\limsup_{s\to\infty}\limsup_{\beta\to\infty}
 \left|h_i^B(mL^s)-h_i^F(mL^s)\right|=0,
 \label{eq:R18}
\end{equation}
where $m$ tends to infinity through sufficiently large dyadic integers.
The inner limit is taken with $m$ and $s$ fixed.
\end{proposition}

\begin{proof}
We use the finite-volume coefficient method of
\cite[Section ``Determining the kinetic drift,'' subsection
``Computing the same coefficients after blocking'']{SharpO4}. Its inputs
here are the finite Taylor maps established in \eqref{eq:R6}, the
Brownian initialization \eqref{eq:R12}--\eqref{eq:R13}, and the matched
coefficient bounds \eqref{eq:R15}. The partition comparison
\eqref{eq:R1}, which concerns center seams on long rectangles, is not
being differentiated.

Fix $m,s$ and put $n=mL^s$. For sufficiently large $\beta$, the
trajectory has at least $s$ layers. After $s$ observation steps the period
is $m$, and the depth index of the original trajectory is $N-s$. We will prove
\begin{equation}
 h_i^M(mL^s)
   =B_{i,N-s}^{q,M}+\mathcal C_i^M(m,s;\beta)
       +O(C_Lm^D e^{-cm})+o_{\beta;m,s}(1),
 \qquad M\in\{F,B\}.
 \label{eq:matching-helicity-decomposition}
\end{equation}
Here $\mathcal C_i^M(m,s;\beta)$ is the Gaussian correction from the
normalized order-two logarithmic expansion at period $m$, computed from
the current canonical coefficients, shapes, and full free kernel. It
includes the contribution of two order-one vertices. Its dependence on
$\beta$ is through those current data; we do not assume that they
converge as $\beta\to\infty$. The leading coupling must be identified
exactly. The constant-order correction and the displayed winding error
will then be controlled uniformly in $s$, whereas the last remainder
needs only tend to zero for fixed $m,s$.

Insert $s$ normalized observation kernels at the actual parameters of
the matched trajectories. Choose the observation blocks inside the fundamental
cell. Near an aligned configuration use the same equivariant formulas as in the renormalization construction; away
from this neighborhood complete them by any local equivariant probability
kernels. These completions need not depend on $\zeta$. For $B$, observe in
the block lift near alignment and then reduce to the quotient. For $F$,
first omit nonzero circle winding and work in its real lift. At fixed
$n$, these omitted sectors and their first two quaternion-twist derivatives
have relative size $O_n(\beta^C e^{-c_n\beta})$: harmonic circle winding
has a positive cost of order $\beta$, and differentiating the finitely
many twisted bonds costs only a polynomial in $\beta$. Thus their
contribution to the helicity is $o(1)$. The normalized observation
insertions leave the retained partition functions unchanged.

Right twisting preserves the common left action on the quaternion field,
so one may pin a last-layer quaternion spin. The common circle translation
may also be pinned. More explicitly, subtract one output circle coordinate
from every lifted coordinate. Integration of the remaining common
translation over one period contributes exactly that period, independently
of the other coordinates. This remains true if it was originally one
microscopic circle value that was required to lie in a fundamental domain.
The extracted phase is identically one on zero-winding configurations,
since $\sum d_t v=0$; it therefore contributes no helicity term.

We next justify the fixed-volume expansion and its identification with the
bulk Taylor maps. After the two symmetries have been removed, the aligned
saddle at $\zeta=0$ is unique and its transverse quadratic form is strictly
positive. Count the angle $\zeta$ at the same scale as the chart fields,
namely $g\asymp\beta^{-1/2}$. At fixed $n,s$ the Wilson expansion is an
ordinary Laplace expansion. For the Brownian expansion, the endpoint and
path estimates used in \eqref{eq:R12}--\eqref{eq:R13} give the same
conclusion. Indeed, restrict endpoint bond distances to $g$ times a fixed
power of $\log(g^{-1})$, and restrict the path-size variables of the
initialization to a corresponding logarithmic bound. The endpoint and
bridge tail estimates, together with the vertical endpoint-bond energy
bound, make the complement smaller than any prescribed power of $g$.
They remain valid after two derivatives in the scaled twist, by the
endpoint-translation estimate. Integer lifts in the horizontal heat
kernel are summed using its Gaussian bounds. Observation kernels similarly
localize their variables in the aligned charts. Thus the relevant Taylor
coefficients are Gaussian polynomial moments. Writing $\eta=\zeta/g$,
expand the logarithm through normalized order two. At fixed $n,s$ its
remainder and its first two $\eta$ derivatives at zero are $o(g^2)$,
uniformly for the actual shape and coordinate-scale parameters in
their compact bands. Since
\[
 h_i^M(n)=-g^{-2}\left.\partial_\eta^2
                    \log Z_i^M(n;g\eta)\right|_{\eta=0},
\]
this remainder contributes $o(1)$ to the unscaled helicity. For fixed $s$,
all observation widths are comparable to $g$.

To identify these coefficients, complete the Gaussians using full kernels
rather than their numerical truncations. Extend each background field
across the boundary with the prescribed right twist. Ambient linear
kernels commute with this continuation, because right multiplication is
orthogonal on $\mathbb R^4$; their compositions and adjoints fold with the
same transport. If $T_x^q$ is the coarse reference field and
$q_x=\exp(\xi_x)T_x^q$, then $\xi_x$ is ordinary periodic. In fact both
$q_x$ and $T_x^q$ acquire the same right multiplier on crossing a period,
which cancels in $q_x(T_x^q)^{-1}$. The affine row in the chart expansion
therefore has the usual folded matrix on $\xi$ and a background term that
is also ordinary periodic. The leading conditional Gaussian precision in
the joint expansion, with both the twist and chart fields counted at scale
$g$, is therefore the ordinary folded precision. Higher-order background
and twist dependence remains among the Taylor vertices. This argument also
applies to horizontal seam strips in \eqref{eq:R12}--\eqref{eq:R13}:
endpoint transport leaves
the left tangent logarithms and bridge variables periodic.

Consequently every Wick contraction without a winding obstruction has the
bulk formula, evaluated on the fields lifted along its contraction paths.
Removing the anchored quaternion rotation cancels the right transports.
The higher-order determinant and Jacobian terms remain among the Taylor
vertices. The logarithmic determinant coefficients are sums over closed
contraction paths and obey the same identification. Repeated transports
on a path contribute at most a fixed power of its length to the twist jets
under consideration, so the exponential path norms still sum absolutely.
Linear circle terms telescope in this saddle calculation.

There is one leading-order point for which an exponentially small estimate
would be insufficient. At harmonic order the first jet of a twisted field
is an affine slope plus an ordinary periodic field. The full free kernels
preserve affine fields. For every folded free edge matrix, translation
invariance makes the cross term between a constant slope and periodic
gradients exactly zero. One can see this also at the last layer by using a
linear ramp that vanishes at the pinned site: its sum with an ordinary
periodic field describes every first jet with the given seam increment.
The affine Hessian normalization therefore gives, exactly, the current
directional coupling $B_i^q$ as the leading unscaled helicity on the square.
The leading determinant is independent of the angle. The strip kernels
have this same property by \eqref{eq:R3}. In particular no term of the form
$\beta$ times an exponentially small winding error is discarded. The
coupling resets in \eqref{eq:R6}, including those at initialization, allow
us to use the actual running couplings: at fixed $s$ the remaining
coefficient error is $o(1)$ and multiplies a chart quadratic.

It remains to control the constant-order terms uniformly as the number of
formal steps grows. At an output period $m'=mL^j$, a discrepancy between a
finite-period Wick coefficient and its folded bulk value requires
contraction paths to travel a distance at least $cm'$. The path estimates
proved for both initializations give a bound $C_Le^{-cm'}$ per anchor;
twist derivatives introduce only fixed path powers. A previously generated
discrepancy grows by at most $C_L$ in one subsequent formal step. This may
be estimated in the total absolute polynomial coefficient norm, including
the number of anchors at the layer where it was generated: prediction and
Wick substitutions have bounded row moments, and a new vertex attached to
an old discrepancy is summed with the same intervertex covariance bound.
The total angle is another field variable in this calculation. Its
occurrences in predictions cost path moments rather than the period.
Products of discrepancies satisfy the same polynomial estimates.

Thus all winding contributions through normalized order two, including
changes of coupling convention, are bounded by a fixed polynomial times
\[
 \sum_{j\ge0}C_L^{j+1}e^{-cmL^j}.
\]
At the final period the pinned Gaussian moments grow at most polynomially
in $m$, by uniform ellipticity. Enlarging a fixed exponent $D$ gives the
bound $C_Lm^D e^{-cm}$ for their contribution to the constant-order
helicity, uniformly in $s$ after the inner limit. Normalized order-one
polynomials have odd total degree in the scaled angle and Gaussian
fields, and hence have zero direct second derivative after Gaussian
integration. Their propagation into order two is already included in the
bound. The estimate includes winding paths whose contractions leave a
function of only a few sites.

This proves \eqref{eq:matching-helicity-decomposition}: earlier bulk
scalar terms have no angle dependence, and the remaining final-period
Gaussian correction is $\mathcal C_i^M(m,s;\beta)$. After separating the
leading coupling, the normalized coefficient slots have fixed caps and
the shapes lie in compact bands. At fixed $m$, pinning removes the kernel
of the uniformly elliptic free precision. The Gaussian moments and
convergent path sums defining the order-two logarithmic correction are
therefore Lipschitz in these data, with a constant depending on $m$ but
independent of $s$.

By \eqref{eq:R15}, the directional coupling discrepancies and canonical
coefficient discrepancies at this layer are $O(\rho^s)$. The free-history
estimates add a term tending to zero with $s$. Subtract
\eqref{eq:matching-helicity-decomposition} for $F$ and $B$, first let
$\beta\to\infty$ at fixed $m,s$, and then let $s\to\infty$. The result
is bounded by $C_Lm^D e^{-cm}$, which tends to zero with $m$. This proves
\eqref{eq:R18}. The use of fixed-volume expansions is legitimate even
when $n$ is smaller than a numerical cutoff required for the exact
retained-density representation: only the full-kernel formal coefficients
are used in this argument.
\end{proof}

\subsection{The microscopic helicities}

We now compute the constant terms that Proposition~\ref{prop:matching-helicity}
compares. Let
\[
 D_j(k)=4\sin^2(k_j/2),\qquad
 D(k)=D_x(k)+D_t(k),\qquad D_*(k)=k_x^2+D_t(k).
\]
For the square of side $n$, write $\int_{W,n}$ for the normalized
momentum sum $n^{-2}\sum_k$ with both momenta in one period of
$(2\pi/n)\mathbb Z$, and write $\int_{B,n}$ for the corresponding sum
with $k_t$ in one period and $k_x\in(2\pi/n)\mathbb Z$ unrestricted.
In both cases the single zero mode $(0,0)$ is omitted.

\begin{lemma}[Gaussian helicity formulas]
\label{lem:matching-microscopic-helicity}
For each fixed even $n$, along the matched parameters as
$\beta\to\infty$,
\begin{equation}
 h_i^F(n)=\beta-\frac32\int_{W,n}\frac{D_i}{D}
                -4\int_{W,n}\frac{\sin^2 k_i}{D^2}+o(1),
 \qquad i=x,t,
 \label{eq:R19}
\end{equation}
and
\begin{equation}
 \begin{aligned}
 h_x^B(n)&=\frac br-4\int_{B,n}\frac{k_x^2}{D_*^2}+o(1),\\
 h_t^B(n)&=br-2\int_{B,n}\frac{D_t}{D_*}
       -4\int_{B,n}\frac{\sin^2 k_t-D_t^2/4}{D_*^2}+o(1).
 \end{aligned}
 \label{eq:R20}
\end{equation}
The errors tend to zero with $n$ fixed. In the correction terms the shape
parameter has already been replaced by its limit $r=1$.
\end{lemma}

\begin{proof}
Distribute the twist uniformly over the links in its direction. We may
use left multiplication, by the inversion symmetry noted above. In the
Wilson model write the imaginary part of a spin, relative to its aligned
value, as $\beta^{-1/2}\pi$, where $\pi$ has three real components. The
leading Gaussian covariance of each component is $D^{-1}$. The second
derivative of the bond factor contributes its scalar product, whose
expectation gives
\[
 \beta-\frac{3}{2}\int_{W,n}\frac{D_i}{D}.
\]
The first derivative contains the chosen component of
$\operatorname{Im}(q_{x+e_i}q_x^{-1})$. Its linear term telescopes on
summing over links. Its quadratic term, up to an irrelevant sign, is
$\pi_x\times\pi_{x+e_i}$. Choosing the first current component, put
\[
 S_i=\sum_x
       \bigl(\pi_x^2\pi_{x+e_i}^3-\pi_x^3\pi_{x+e_i}^2\bigr)
     =\sum_x\pi_x^2(\pi_{x+e_i}^3-\pi_{x-e_i}^3).
\]
Wick contraction yields
\[
 n^{-2}\mathbb E S_i^2
       =4\int_{W,n}\frac{\sin^2 k_i}{D^2}.
\]
Subtracting this current variance proves \eqref{eq:R19}. Every quantity
in this calculation is unchanged by adding a common field constant.
Thus pinning one spin or using the mean-zero propagator gives the same
answer. Nonzero circle windings in $F$ are exponentially small at fixed
$n$ and do not alter this formula.

For $B$, the Brownian quadratic energy and \eqref{eq:R13} give four
independent tangent fields at scale $b^{-1/2}$: three quaternion
components $\pi$ and the circle component, denoted $\nu$. At $r=1$ their
covariance is $D_*^{-1}$. The vertical scalar product is multiplied by
$\exp(i\Delta_t v)$, so all four fields contribute to its quadratic
correction. This gives $br-2\int_{B,n}D_t/D_*$. The linear circle term
again telescopes. The quadratic vertical current now has the form
$S_t\pm iT_t$, where the horizontal coordinate is continuous:
\[
 \begin{aligned}
 S_t&=\sum_{u=0}^{n-1}\int_0^n
    \bigl(\pi^2(x,u)\pi^3(x,u+1)
          -\pi^3(x,u)\pi^2(x,u+1)\bigr)\,dx,\\
 T_t&=\sum_{u=0}^{n-1}\int_0^n
         (\Delta_t\nu)(x,u)(\Delta_t\pi^1)(x,u)\,dx.
 \end{aligned}
\]
The components occurring in $S_t$ and $T_t$ are independent. Hence
\[
 n^{-2}\mathbb E S_tT_t=0,\qquad
 n^{-2}\mathbb E T_t^2=\int_{B,n}\frac{D_t^2}{D_*^2}.
\]
The derivative of the logarithm of a complex weight uses the ordinary
bilinear second moment, with no complex conjugation. Thus the $iT_t$
term contributes the negative of this last quantity to the current
variance. Together with the Wilson current calculation, this is exactly
the vertical formula in \eqref{eq:R20}.

For the horizontal response, rotate each Brownian path steadily on the
left. The deterministic-rotation identity for the heat measure used in
\eqref{eq:R12}, in Stratonovich convention, gives the constant quadratic
stiffness $b/r$. Its first derivative is the integrated current
$\int\operatorname{Im}(dq\,q^{-1})$ in direction $e$, multiplied by
$b/r$ and by the angle-per-period factor. The linear tangent term
telescopes. The remaining Gaussian current is
\[
 \sum_t\int_0^n(\pi^2\,d\pi^3-\pi^3\,d\pi^2),
\]
whose variance, divided by $n^2$, is
$4\int_{B,n}k_x^2/D_*^2$.

For completeness, this last identity can be obtained without treating a
Brownian derivative as a pointwise field. Interpolate the paths on a
dyadic polygonal mesh of spacing $\ell=n/N'$, including all strip
endpoints, and let $N'\to\infty$. The leading fields consist of
independent Gaussian bridges with the deterministic interpolation and
Gaussian $C^1$ additions from the strip completion. Polygonal
Stratonovich approximation therefore converges with every fixed moment;
there are no jumps at the finitely many joins. The point covariance has
an absolutely summable horizontal Fourier series. Sampling replaces its
coefficient at $-N'/2<j\le N'/2$ by the sum of coefficients at
$j+uN'$, $u\in\mathbb Z$, with the zero mode still omitted. The
antisymmetric polygonal area has Wick multiplier
\[
 \frac{4\sin^2(2\pi j/N')}{\ell^2}.
\]
At fixed $n$, the added covariance aliases are $O_n(N'^{-2})$ uniformly
in $j$. The multiplier is $O_n(j^2)$ and the original covariance is
$O_n((1+|j|)^{-2})$, so summing the alias errors gives a quantity tending
to zero. The remaining sums converge to the displayed Fourier variance.
Subtracting a constant for pinning does not change the area. This proves
the horizontal formula. Finally, since $r\to1$, continuity of each
fixed-$n$ correction permits replacing $r$ by $1$ there, while retaining
$b/r$ and $br$ in the leading terms.
\end{proof}

\begin{proposition}[The finite coupling shift]
\label{prop:matching-coupling-shift}
For every choice of matched parameters satisfying \eqref{eq:R15},
\begin{equation}
 \beta-b=\frac14-\frac{1+\log 2}{2\pi}+o(1)
 \qquad (\beta\to\infty).
 \label{eq:R21}
\end{equation}
\end{proposition}

\begin{proof}
Average the two directional formulas in
Lemma~\ref{lem:matching-microscopic-helicity} and compare them by
\eqref{eq:R18}. The leading Brownian average is
\[
 \frac b2(r+r^{-1})=b+o(1),
\]
since $r-1=O(\beta^{-1})$. The identities
$\sin^2 k_i=D_i-D_i^2/4$ express the difference of the constant terms as
paired momentum sums. We must pair them before passing to integrals,
since each propagator sum diverges as $n\to\infty$. On their common
momentum region,
\[
 \frac1{D_*}-\frac1D
   =\frac{D_x-k_x^2}{D_*D}=O(1)
 \qquad(k\to0).
\]
The other resulting integrands are locally bounded, and the Brownian
tails in $k_x$ are absolutely and uniformly summable. Riemann convergence
therefore applies to the paired expression.

To respect the order in \eqref{eq:R18}, start with any subsequence along
which the bounded difference $\beta-b$ converges. First take its
$\beta$ limit at fixed $m,s$, then let $s$ and $m$ tend to infinity as
in \eqref{eq:R18}. This determines that limit point from the paired
integrals. Equivalently,
\begin{equation}
 \begin{split}
 \beta-b={}&\frac34-\int_B\frac{D_t}{D_*}
                  +\int_B\frac{D_t^2}{D_*^2}
           -\frac12\int_W\frac{D_x^2+D_t^2}{D^2}\\
 &\hspace{12mm}-2\left(\int_B\frac1{D_*}-\int_W\frac1D\right)+o(1).
 \end{split}
 \label{eq:matching-integral-difference}
\end{equation}
Here $\int_W$ means $(2\pi)^{-2}$ times integration over
$[-\pi,\pi]^2$, and $\int_B$ means the same normalization over
$\mathbb R\times[-\pi,\pi]$. The propagator difference is taken with a
common neighborhood of the origin deleted before its size tends to zero.
In the finite sums the constant $3/4$ is replaced by
$3(1-n^{-2})/4$, which has the same limit. It remains to evaluate the
constants in \eqref{eq:matching-integral-difference}.

Set $d=2|\sin(k_t/2)|$ and use a vertical angular average denoted by
$\langle\cdot\rangle$. Horizontal integration gives
\[
 \int_{\mathbb R}\frac{dk_x}{2\pi(k_x^2+d^2)}=\frac1{2d},
 \qquad
 \int_{\mathbb R}\frac{dk_x}{2\pi(k_x^2+d^2)^2}=\frac1{4d^3}.
\]
Consequently
\[
 \int_B\frac{D_t}{D_*}=\langle d/2\rangle=\frac2\pi,
 \qquad
 \int_B\frac{D_t^2}{D_*^2}=\langle d/4\rangle=\frac1\pi.
\]
By symmetry, the lattice term with coefficient $1/2$ equals
$\int_W D_t^2/D^2$. The one-dimensional lattice propagator and its
derivative yield
\[
 \int_W\frac{D_t^2}{D^2}
   =\left\langle\frac{d(d^2+2)}{(d^2+4)^{3/2}}\right\rangle
   =\frac1\pi\int_0^1\frac{3-2u^2}{(2-u^2)^{3/2}}\,du
   =\frac12-\frac1{2\pi}.
\]
For the last equality an antiderivative is
$2\arcsin(u/\sqrt2)-u/(2\sqrt{2-u^2})$.

The propagator difference is
\[
 \left\langle\frac1{2d}-\frac1{d\sqrt{d^2+4}}\right\rangle.
\]
To evaluate it, define
\[
 I(z)=\int_0^{\pi/2}
     \frac{1-(1+z\sin^2 y)^{-1/2}}{\sin y}\,dy,
 \qquad z\ge0.
\]
Differentiation under the integral gives
$I'(z)=1/(2(1+z))$, and $I(0)=0$. The required difference is therefore
$I(1)/(2\pi)=\log2/(4\pi)$. Substitution in
\eqref{eq:matching-integral-difference} gives
\[
 \frac34-\frac2\pi+\frac1\pi
       -\left(\frac12-\frac1{2\pi}\right)-\frac{\log2}{2\pi}
   =\frac14-\frac{1+\log2}{2\pi}.
\]
Every limit point is this number, proving \eqref{eq:R21}.
\end{proof}

\section{Recovery of the full lattice mass}
\label{sec:assembly}

We now apply the comparison to rectangles whose odd traces are small but
still much larger than the comparison error. The additional circle field
has a massless Gaussian contribution; time doubling removes its extensive
part and leaves an error small enough to recover the Brownian transfer
norms. A nearby-fugacity estimate then verifies the particle-number
hypothesis of \eqref{eq:B1}.

Fix the admitted depth $N$ chosen for the Wilson coupling $\beta$, and put
$a=L^{-N}$. Choose dyadic terminal lengths
\begin{equation}
 w\asymp\beta(\log\beta)^2,
 \qquad t\asymp\sqrt\beta(\log\beta)^4,
 \qquad X=\frac wa,\quad n=\frac ta,\quad W=\frac{rX}{b}.
 \label{eq:assembly-rectangles}
\end{equation}
All fine periods are even integers, and we also use time length $2t$.
The short terminal period tends to infinity and hence eventually exceeds
the fixed admission, stencil, and tile sizes. The comparisons needed below
are summarized by
\begin{equation}
 t=o(\beta),\qquad t=o(w),\qquad
 \frac{\beta w}{t}\gg t,\qquad
 \frac{\beta t}{w}\longrightarrow\infty,
 \qquad \frac{w}{t}=o(t).
 \label{eq:assembly-scale-separations}
\end{equation}
For example, $(\beta w/t)/t\asymp\beta/(\log\beta)^6$ and
$w/t^2\asymp(\log\beta)^{-6}$.

\subsection{The Wilson parity traces}

Write $Z^q_{\varepsilon,\eta}(t,w)$ for the Wilson partition function,
where $\varepsilon$ and $\eta$ specify the time and space center seams:
$+$ is periodic and $-$ is the minus-center twist. With periodic space,
define
\[
 Z_l^q(t)=\frac12\bigl(Z^q_{+,+}(t,w)+(-1)^lZ^q_{-,+}(t,w)\bigr),
 \qquad l=0,1.
\]
These are the even and odd row-transfer traces, because the time-center
seam inserts total spin inversion. Define $Z_l^{F(\phi)}$ and $Z_l^B$
by the same time-seam half-sum and half-difference.

\begin{lemma}[Wilson trace asymptotics]
\label{lem:assembly-wilson-traces}
For the rectangles \eqref{eq:assembly-rectangles}, uniformly along the
matched trajectories,
\begin{equation}
 \log\frac{Z_1^q(t)}{Z_0^q(t)}
       =-\frac{t}{a}m_{\rm lat}(\beta)+o(t),
 \qquad
 2\log Z_0^q(t)-\log Z_0^q(2t)=o(1).
 \label{eq:C1}
\end{equation}
The first formula also holds with $t$ replaced by $2t$.
\end{lemma}

\begin{proof}
Let $U(w)$ be the row transfer normalized by its top eigenvalue, and let
$E(w)=-a^{-1}\log\|U(w)|_{\Omega^\perp}\|$. Proposition~\ref{prop:cylinder-rate} gives
\[
 E(w)=\frac{m_{\rm lat}(\beta)}a+O(w^{-1/4}),
 \qquad \frac{m_{\rm lat}(\beta)}a\asymp1.
\]
By Lemma~\ref{lem:wilson-detector}, the largest nonvacuum eigenvalue
is attained in odd parity. Write
$s(u,w)=\operatorname{Tr}U(w)^{u/a}-1$. The long-rectangle estimate
\eqref{eq:wilson-long-rectangle}, together with
\eqref{eq:rectangle-identities}, gives exponentially small trace tails at
$u=t,2t$, and also at a dyadic $t_0\asymp t^\epsilon$ for any fixed
sufficiently small $\epsilon>0$. Indeed their prefactors grow at most
polynomially in the present $w,t$.

After dividing by the top eigenvalue to the power $n$, the odd trace is
at least $e^{-tE(w)}$ and at most
\[
 s(t_0,w)e^{-(t-t_0)E(w)}\le e^{-(t-t_0)E(w)}.
\]
Its logarithm is therefore $-tE(w)+o(t)$. The normalized even trace is
$1+o(1)$, since it contains the vacuum and its remaining terms are bounded
by $s(t,w)$. These estimates, also at $2t$, prove \eqref{eq:C1}; the
vacuum eigenvalue cancels in the time-doubling expression.
\end{proof}

\subsection{Removing the extra circle field}

Let $Z_G(t)$ be the zero-winding Gaussian circle integral in $F$ on the
$n\times X$ rectangle. Its value is independent of the quaternion field.
The phase $\exp(i\phi\sum d_t v)$ is one in this sector.

\begin{lemma}[The circle contribution]
\label{lem:assembly-circle}
For $l=0,1$, at both time lengths $t$ and $2t$,
\begin{equation}
 Z_l^{F(\phi)}(t)=Z_G(t)Z_l^q(t)(1+o(1)).
 \label{eq:C2}
\end{equation}
Moreover,
\begin{equation}
 2\log Z_G(t)-\log Z_G(2t)=o(t).
 \label{eq:C3}
\end{equation}
\end{lemma}

\begin{proof}
The circle integral is the product of its zero-winding Gaussian integral
and a sum of harmonic winding weights. A temporal winding has Gaussian
cost at least $c\beta w/t$ times its winding number squared; a spatial
winding has cost at least $c\beta t/w$ times its winding number squared.
Each winding also inserts the corresponding center signs in the
quaternion boundary conditions, and the temporal winding may carry the
phase of modulus one.

We first estimate the effect of a spatial quaternion twist. Expanding
all Wilson bonds in powers gives nonnegative integrated coefficients
before seam signs are inserted. Consequently, for either time seam,
\[
 \bigl|Z^q_{\varepsilon,+}-Z^q_{\varepsilon,-}\bigr|
      \le Z^q_{+,+}-Z^q_{+,-}.
\]
Applying the row-trace estimate along the other lattice axis bounds the
right-hand side by $Ce^{-cw}Z^q_{+,+}$. This also bounds the change in
either time-parity combination.

By \eqref{eq:C1}, $Z_1^q/Z_0^q\ge e^{-Ct}$ for a fixed $C$ and all
large $\beta$. The total contribution of nonzero temporal winding is
therefore negligible even relative to the odd trace, since
\[
 e^{Ct-c\beta w/t}\longrightarrow0.
\]
For spatial winding we use the preceding twist estimate instead of
requiring its Gaussian cost to dominate $t$. The sum of its nonzero
relative Gaussian weights tends to zero because $\beta t/w\to\infty$.
Replacing each spatially twisted quaternion partition by the untwisted
one costs, relative to either parity trace, at most a constant times
$e^{Ct-cw}$, which also tends to zero. Summing the Gaussian winding tails
proves \eqref{eq:C2}. The same comparisons hold at $2t$. This argument
controls the odd half-difference itself, rather than subtracting two
partition estimates whose errors might exceed that difference.

To prove \eqref{eq:C3}, use the initial circle variable of period
$\ell_v=\pi$ and its kinetic coefficients $c_x,c_t\asymp\beta$; their
ratio is bounded above and below. Up to a scalar exponential proportional
to $nX$, the zero-winding Gaussian integral is
\[
 Z_G(t)=
 \frac{\ell_v\sqrt{nX}(2\pi/c_t)^{(nX-1)/2}}
      {\sqrt{\det{}'(D_t+(c_x/c_t)D_x)}}.
\]
The prime omits the constant mode. Factor the determinant over horizontal
momenta $p_x\in(2\pi/X)\mathbb Z$ in one period. With
\[
 2\cosh\gamma(p_x)=2+(c_x/c_t)D_x(p_x),
\]
the elementary product over time momenta gives
\[
 \sqrt{\det{}'(D_t+(c_x/c_t)D_x)}
       =n\prod_{p_x\ne0}2\sinh\bigl(n\gamma(p_x)/2\bigr).
\]
In $2\log Z_G(t)-\log Z_G(2t)$ all terms linear in $n$ cancel.
The product terms $\log(1-e^{-n\gamma(p_x)})$ have total absolute size
$O(1+X/n)$: use $\gamma(p_x)\ge c|p_x|$ for centered momenta and the
integrability of $|\log(1-e^{-u})|$ on $(0,\infty)$. The remaining
$n,X$ prefactor is $\sqrt{X/n}$, so it contributes only a logarithm of
the aspect ratio, together with $O(\log\beta)$. Since $X/n=w/t=o(t)$,
these terms prove \eqref{eq:C3}.
\end{proof}

\subsection{From traces to the Brownian transfer norms}

The comparison \eqref{eq:R1} has an absolute error $e^{-c\beta}$ relative
to the untwisted flat reference. Because $t=o(\beta)$, this is negligible
relative to the odd expression in \eqref{eq:C2}. The volume scalars in
\eqref{eq:R1} are independent of the center seam and proportional to the
rectangle area. They cancel both in parity ratios and in time doubling.
Combining \eqref{eq:C1}--\eqref{eq:C3} therefore gives
\begin{equation}
 \log\frac{Z_1^B(t)}{Z_0^B(t)}
      =-\frac{t}{a}m_{\rm lat}(\beta)+o(t),
 \qquad
 2\log Z_l^B(t)-\log Z_l^B(2t)=o(t),\quad l=0,1.
 \label{eq:C4}
\end{equation}
For $l=1$, add twice the log-ratio relation at $t$ minus the same
relation at $2t$ to the even-parity doubling relation. The terms linear
in the mass cancel, giving the asserted odd-parity relation.

We record the elementary spectral fact that makes these estimates useful.
It allows the number of appreciable eigenvalues to depend on the
parameters.

\begin{lemma}[Trace doubling and the norm]
\label{lem:assembly-trace-norm}
Let $A$ be a nonzero positive compact operator. Suppose $A^n$ is trace
class for a positive integer $n$, and write $\tau=\|A\|$. Then
\[
 0\le\log\frac{\operatorname{Tr}A^n}{\tau^n}
 \le 2\log\operatorname{Tr}A^n-\log\operatorname{Tr}A^{2n}.
\]
The same conclusion holds for a direct sum of positive operators whenever
the indicated total trace is finite.
\end{lemma}

\begin{proof}
Let $x_j$ be the eigenvalues of $A/\tau$, with multiplicity. They belong
to $[0,1]$, and at least one equals $1$. Thus
$1\le\sum_jx_j^{2n}\le\sum_jx_j^n$. Taking logarithms gives the
claim. For a direct sum, finiteness of the positive total trace makes it
compact and ensures that its nonzero norm is attained; the same argument
applies.
\end{proof}

The Brownian operators constructed in Section~\ref{sec:brownian} are
positive and have finite total trace in each parity. Let $\tau_l$ be the
maximum of their norms over particle numbers of parity $l$. Both parities
have nonzero trace. Lemma~\ref{lem:assembly-trace-norm} and
\eqref{eq:C4} give
\[
 \log Z_l^B(t)=n\log\tau_l+o(t),\qquad l=0,1.
\]
Subtracting and using $n=t/a$ yields
\begin{equation}
 \log\tau_0-\log\tau_1=m_{\rm lat}(\beta)+o(a).
 \label{eq:C5}
\end{equation}
The error is $o(a)$ because the preceding errors were $o(t)$ before
division by $n$. This is the precision required for the exponentially
small lattice mass.

\begin{lemma}[Location of maximizing particle numbers]
\label{lem:assembly-particle-density}
For either parity, every particle number $j$ attaining $\tau_l$ on the
width \eqref{eq:assembly-rectangles} satisfies
\[
 \frac jW=b^2+O(b),
\]
with a constant uniform as $\beta\to\infty$.
\end{lemma}

\begin{proof}
Change the fugacity by $\exp(\pm1/b)$. In the notation of
\eqref{eq:R17}, choose $b',r'$ so that
\[
 (b'/b)^2=e^{\pm1/b},\qquad r'/b'=r/b.
\]
The heat time and the width remain fixed. Homogeneity in fugacity
multiplies the sector-$j$ transfer by $e^{\pm j/b}$. The numerator
untwisted trace in \eqref{eq:R17} is therefore at least
$\tau_l^n e^{\pm nj/b}$. On the other hand,
\eqref{eq:C4}--\eqref{eq:C5} bound the total trace at the matched
parameter by $\tau_l^n e^{Ct}$ for either $l$. Hence
\[
 \pm\frac{nj}{b}
   \le Xn\bigl(b'r'-br+O(b^{-1})\bigr)+O(t).
\]
Since $b'r'=br e^{\pm1/b}$, division by $n$ gives, for the two signs,
\[
 \begin{aligned}
 j/b&\le X\bigl(r+O(b^{-1})\bigr)+O(a),\\
 -j/b&\le X\bigl(-r+O(b^{-1})\bigr)+O(a).
 \end{aligned}
\]
Using $X=Wb/r$ and $r\asymp1$, these are the required upper and lower
bounds $j/W=b^2+O(b)$. The argument uses only the pressure estimate
\eqref{eq:R17} at the displaced fugacity; it requires no matched
trajectory at that parameter.
\end{proof}

\begin{proof}[Proof of Theorem~\ref{thm:main}]
The depth calibration gives
$a\asymp\sqrt\beta e^{-\pi\beta}$. Since $b=\beta+O(1)$ and
$r=1+O(\beta^{-1})$, the width in \eqref{eq:assembly-rectangles}
satisfies
\[
 W=\frac{rw}{ab}
   \asymp e^{\pi b}b^{-1/2}(\log b)^2.
\]
Thus the width hypothesis of \eqref{eq:B1} holds, with logarithmic
exponent $2$. Lemma~\ref{lem:assembly-particle-density} verifies its
particle-number hypothesis for every maximizing sector in both parities.
We may therefore apply \eqref{eq:B1} to \eqref{eq:C5}, obtaining
\[
 m_{\rm lat}(\beta)
    =32\sqrt2\,b^{1/2}e^{-\pi b}(1+o(1))+o(a).
\]
Multiplication by $e^{\pi\beta}/\sqrt\beta$ turns the last error into
$o(1)$, by the depth calibration. Proposition~\ref{prop:matching-coupling-shift}
gives the remaining factor:
\[
 \begin{split}
 \lim_{\beta\to\infty}
       \frac{e^{\pi\beta}}{\sqrt\beta}m_{\rm lat}(\beta)
  &=32\sqrt2\,
       \exp\!\left(\frac\pi4-\frac{1+\log2}{2}\right)\\
  &=32\exp\!\left(\frac\pi4-\frac12\right).
 \end{split}
\]
The quantity entering \eqref{eq:C1} and \eqref{eq:C5} is the full
Osterwalder--Schrader transfer gap identified in
Sections~\ref{sec:sector}--\ref{sec:cylinders}, including the
rotation-invariant sectors. The asserted asymptotic therefore has the
full-gap meaning of Theorem~\ref{thm:main}.
\end{proof}
